\documentclass[10pt]{article}
\usepackage[margin=20mm]{geometry}
\usepackage{amsmath,amssymb,mathtools,fontspec,unicode-math,xurl,hyperref}
\IfFontExistsTF{DejaVu Serif}{\setmainfont{DejaVu Serif}}{\setmainfont{Latin Modern Roman}}
\setmathfont{latinmodern-math.otf}
\setlength{\parindent}{0pt}\setlength{\parskip}{6pt}\sloppy
\title{Technical review of The \(L_p\)-Minkowski problem with super-critical exponents}
\author{AI-assisted technical review}
\date{}
\begin{document}
\maketitle
Abstract: The existence argument is examined in the bound publisher version. An independent matrix parametrization verifies the top homology of the central ellipsoid space, and the energy identity, stopped-flow selection, and stationary subsequence argument are checked. No substantive error is established within this scope. The detailed stratified homology calculation and every curvature estimate are not individually certified, and actual page-image inspection remains incomplete.\par
\section*{Identity and review scope}
The paper is by Qiang Guang, Qi-Rui Li, and Xu-Jia Wang, The \(L_p\)-Minkowski problem with super-critical exponents, DOI 10.4171/jems/1733 (https://doi.org/10.4171/jems/1733), J. Eur. Math. Soc. 28 (2026), 735–775 [paper]. The exact reviewed file is the publisher version, first published online on 13 November 2025. The existing 2025 frozen identifier is retained. PDF page \(j\) corresponds to printed page \(734+j\).\par
Actual text reading covers PDF pages 1–20 and 32–36. Targets are Theorem 1.1, the monotonicity and compactness mechanism, Theorem 3.7 and its role in Lemma 3.8, the subsequence limit, and the architecture of the curvature estimates. Pages 21–31 and most of Appendix A were not individually audited. A separate exact argument below checks the central homology statement without claiming to have checked every step of Section 4.\par
The range is \(p<-n-1\), \(n\geq1\). The support function must be positive and arise from a convex body containing the origin in its interior. The smooth theorem assumes positive \(f\in C^{1,1}(S^n)\); the weak corollary retains two positive bounds on the density. The review does not replace these with mere nonnegativity, assume evenness, or assert uniqueness. The publisher metadata confirms the negative signs in the super-critical range; the PDF text extraction drops some mathematical glyphs, so ambiguous extraction alone is not used to diagnose errors.\par
\section*{Checked analytic mechanism}
Write \(D=\det(\nabla^2u+uI)=K^{-1}\). Equation (2.4) is \[u_t=u-\frac{fu^p}{D}.\] The support-function variation of volume gives the independent identity \[\begin{aligned}
\frac{d}{dt}\left(\operatorname{Vol}(\Omega)
      -\frac1p\int_{S^n}fu^p\,d\sigma\right)
 &=\int_{S^n}(D-fu^{p-1})u_t\,d\sigma\\
 &=\int_{S^n}\frac{u}{D}(D-fu^{p-1})^2\,d\sigma\geq0.\end{aligned}\] Under the positive and uniformly convex hypotheses, equality is equivalent to the elliptic equation. This agrees with Lemma 2.1 (PDF page 7, printed page 741).\par
Lemmas 2.2–2.4 address degeneration toward the origin, extreme volume, and large eccentricity. In Lemma 2.2, the volume and polar-cone estimates combine into a lower bound proportional to \(d^{-(-p-n-1)}\) for \(J^2\). The exponent is positive precisely in the stated super-critical range. No corresponding coercive conclusion at \(p=-n-1\) is inferred. Together with the John ellipsoid inclusions, an upper bound for \(J\) gives both inner and outer radius control.\par
The stopped flow in Section 3 is part of the topological construction. It is stationary at energy at least \(9A_0/10\), whereas the selected trajectories at the prescribed times have energy at most \(A_0/2\). The monotonicity and the continuous-dependence argument in Remark 3.11 are needed to ensure the limiting selected trajectory never reaches the stopping threshold. Lemma 3.12 gives the stricter \(3A_0/4\) bound. Thus one cannot infer compactness for arbitrary initial data merely from energy monotonicity; the selected initial data supply the missing upper bound.\par
The curvature-estimate architecture is consistent: an upper bound for the product of the principal curvatures, together with a positive lower bound for each curvature, also bounds each curvature above. Equivalently, the eigenvalues of \(\nabla^2u+uI\) stay between positive constants. Sections 2.2 and 5 then invoke parabolic and elliptic regularity. The auxiliary-function argument and its endpoint cases were read, but every differentiated tensor coefficient and each cited regularity theorem were not independently reconstructed.\par
\section*{Independent verification of the central ellipsoid homology}
Let \(m=n+1\) and let \(\mathcal E\) be the space called \(E\) in (3.4): origin-centred ellipsoids of volume \(\operatorname{Vol}(B_1)\) with eccentricity \(\bar e>1\), carrying the Hausdorff topology. This is the boundary of the corresponding bounded-eccentricity set of symmetric ellipsoids, not the larger space allowing translation.\par
Every such ellipsoid has a unique representation \[E_A=A B_1,\qquad A=A^{\mathsf T}>0,\quad \det A=1.\] The eigenvalues of \(A\) are its semiaxis lengths. This representation and its inverse are continuous: the support function is \(h_{E_A}(x)=|Ax|\), so \(h_{E_A}^2\) determines the symmetric matrix \(A^2\), and the positive square root is unique and continuous. Repeated semiaxes therefore create no choice of eigenbasis in this parametrization.\par
Set \(H=\log A\). This identifies determinant-one positive definite matrices with \[\mathcal V=\{H=H^{\mathsf T}:\operatorname{tr}H=0\},
\qquad \dim\mathcal V=\frac{m(m+1)}2-1.\] For nonzero \(H\in\mathcal V\), define \[w(H)=\lambda_{\max}(H)-\lambda_{\min}(H)>0.\] It is continuous and satisfies \(w(rH)=r\,w(H)\) for \(r>0\). Moreover, \(w(H)=0\) in \(\mathcal V\) forces \(H=0\). The eccentricity condition is exactly \[w(H)=\log\bar e.\] Consequently the map from the Frobenius unit sphere in \(\mathcal V\), \[H\longmapsto \frac{\log\bar e}{w(H)}H,\] is a homeomorphism onto this spectral-width level set; its inverse is radial normalization. Composing with the matrix exponential and \(A\mapsto E_A\) proves \[\mathcal E\cong S^q,\qquad
q=\frac{(n+1)(n+2)}2-2
  =\frac{n(n+1)}2+n-1.\] In particular, \[H_q(\mathcal E;\mathbb Z)=\mathbb Z,\qquad
H_k(\mathcal E;\mathbb Z)=0\quad(k>q).\] This verifies Theorem 3.7 and the higher-degree vanishing used in (4.31). It is an analytic/topological proof, not a finite numerical experiment.\par
In the contradiction on PDF page 17, take \(k=q+n>q\). The Künneth term \(H_{k-n}(\mathcal E)=H_q(\mathcal E)=\mathbb Z\) survives, whereas \(H_k(\mathcal E)\) and \(H_k(S^n)\) vanish. Thus the contradiction does not require cancellation of an unspecified infinite abelian group. The surrounding retraction and Mayer–Vietoris constructions are inspected at the level described in the coverage ledger; the independent sphere model verifies the decisive homology input.\par
\section*{Limit, outcome, and boundaries}
With the selected trajectory uniformly controlled, \(J\) is bounded and nondecreasing. Therefore there exist \(t_i\to\infty\) with \(J'(t_i)\to0\). Compactness and the nonnegative square identity yield a stationary subsequential limit solving (1.1). This suffices for Theorem 1.1. No additional claim of convergence of the entire trajectory to a unique limit is certified by this review. A compactness statement in a chosen Hölder exponent can use the stated estimate at a slightly larger exponent.\par
The weak-density passage in Corollary 1.2 uses bounds independent of smooth derivatives of the approximating \(f_j\), followed by weak Monge–Ampère convergence. Its cited strict-convexity and \(C^{1,\alpha}\) regularity inputs were not reproved here. No supported substantive defect has been found in the checked existence mechanism.\par
Item | State | Scope\par
Issue judgment | No Issue Found | Inspected existence mechanism\par
Independent proof | Verified | Ellipsoid sphere model and top homology\par
Repair/adoption | Not applicable / not requested | No defect requiring a proposed fix\par
Scientific decision | Pending | Initial review, no human acceptance\par
Remaining work | Actual image inspection | Source evidence and every report page\par
Bounded searches and acquired metadata identified no separate correction, supplement, or released computational repository for this work. Absence is not established. No author code was executed and no earlier preprint was substituted for the bound publisher PDF.\par
\section*{Supporting checks}
> 00000TOKEN.\par
Actual commands, exit codes, TeX and compiler hashes, retained per-paper dependency payloads, PDF correspondence, and inventories accompany this report in `verification/0354`. The current assistant performs the scientific reasoning and directs native tools for mechanical operations. Controller UI observations are separately attributed; backend equivalence is not independently verified. No human review is recorded.\par
9 Q. Guang, Q.-R. Li and X.-J. Wang, The \(L_p\)-Minkowski problem with super-critical exponents, J. Eur. Math. Soc. 28 (2026), 735–775, DOI:10.4171/jems/1733, bound publisher PDF. EMS Press, article metadata, https://ems.press/journals/jems/articles/14299304, consulted 27 October 2026.\par
\end{document}
